\section{Introduction}

Container-based workloads have become a cornerstone of modern cloud infrastructure, revolutionizing how applications are developed, deployed, and scaled, particularly within microservices architectures and DevOps workflows. Unlike hypervisor-based virtual machines (VMs), which each require a full operating system and rely on a hypervisor for hardware virtualization, containers share the host system's kernel, reducing their overhead. This efficiency makes them exceptionally well-suited for continuous integration and rapid scaling, allowing organizations to optimize resource utilization and reduce costs while ensuring quick, consistent deployments. Furthermore, containers enable seamless portability by packaging all userspace dependencies into a unified image, thereby avoiding the common problem of environment-specific discrepancies. Yet, despite containers' numerous benefits, they also introduce new security challenges, especially where container confinement is concerned.

Linux kernel features such as cgroups and namespaces provide resource, process, and network isolation for containers. However, this separation is not foolproof. Kernel vulnerabilities, such as buffer overflows or privilege escalation flaws can allow a process to circumvent the very mechanisms designed to confine it. Misconfigurations can similarly open holes in container security, for instance by granting unnecessary capabilities, mounting sensitive host directories, or running containers in privileged mode. Container escapes exploit these weaknesses to break out of the container's resource and namespace isolation. Once an attacker gains elevated privileges or direct access to the host system, they can disrupt other containers, exfiltrate data, or compromise the entire infrastructure. 

Technologies like SELinux, AppArmor, and seccomp provide additional layers of host-based security to mitigate these issues but are notoriously prone to misconfiguration. Best practices guides for container security~\cite{CIS} can help but they offer no guarantee that a given configuration will perform as intended by the end user. As container adoption has skyrocketed across industries, particularly in multi-tenant cloud environments, these security challenges have become more pronounced.

To improve container security, the cloud industry has developed technologies like Kata Containers\footnote{\url{https://katacontainers.io}}, which isolate containers further by running them within lightweight VMs. While this approach can be effective, it also adds another layer of complexity which is subject to its own set of vulnerabilities and implementation errors. Simply adding additional layers of security does not, it itself, mean that security goals have been achieved.

Test suites play a crucial role in verifying that systems meet specified requirements and function correctly under defined conditions. They are widely used across computing disciplines, from performance benchmarking to security assessment. For instance, the Standard Performance Evaluation Corporation (SPEC) \cite{henning2000spec} is best known for developing industry-standard benchmarks that evaluate CPU performance and power efficiency, while the Transaction Processing Performance Council (TPC) \cite{10.1007/978-3-642-18206-8_1} establishes benchmarking standards for database performance.
%Beyond performance testing, test suites are essential in software development to prevent regressions. Many organizations mandate that all code changes pass a comprehensive suite of tests before integration. In production environments, test suites help maintain reliability through uptime monitoring, ensuring that services remain operational. 
Specialized security-focused test suites play a vital role in detecting vulnerabilities in deployed systems. Tools such as network scanners aid in identifying misconfigurations, outdated software versions, and unintended service exposures, helping organizations proactively mitigate security risks. However, a significant gap remains in the context of containers. There are currently no standardized or widely adopted test suites designed specifically to evaluate container confinement, functionality, or isolation effectiveness. This absence leaves organizations to rely on ad hoc methods, which are often inadequate for uncovering subtle misconfigurations or ensuring robust confinement. Given the growing reliance on containerized environments in modern infrastructure, the development of comprehensive test suites tailored for containers is imperative.

In this paper, we describe \houdini, the first test suite for verifying container confinement. Given a Docker container configuration (including the kernel, Docker, and runc versions), \houdini will instantiate the container in a standalone QEMU-based VM and perform multiple tests (or \textit{tricks}, as we call them) to evaluate the container's security posture. These tests explore vulnerabilities such as improper resource allocation, inadequate isolation, and insecure filesystem setups. \houdini is written in Python and provides an extension language for writing tricks, allowing it to adapt as new container escape vulnerabilities arise.

%In this paper, we describe the motivation, design, and implementation of Houdini. We also present the results of case studies by showing how Houdini can be used to detect misconfigured containers. 

%maybe better for the conclusion
%Our hope with Houdini is that it will help with the deployment of better confined containers and will support the development of new technologies to more reliably confine containers.

This paper makes the following contributions:
\begin{itemize}
    \item Defines container confinement and its associated challenges.
    \item Introduces \houdini, a modular framework for testing container security.
    \item Demonstrates Houdini's ability to detect security issues, misconfigurations, and inconsistencies in container confinement.
\end{itemize}

The rest of this paper proceeds as follows. Section~\ref{sec:container-confinment-problem} defines container confinement and associated challenges. Section~\ref{sec:motivation} further motivates the need for a container security test suite. We describe \houdini's design and implementation in Sections~\ref{sec:design} and~\ref{sec:implementation}. Section~\ref{sec:evaluation} presents three case studies demonstrating how \houdini can detect security issues, misconfigurations, and other inconsistencies. We review related work in Section~\ref{sec:related} and conclude in Section~\ref{sec:conclusion}.
